\documentclass[12pt,a4paper]{article}
%\usepackage[in]{fullpage}
\usepackage[utf8x]{inputenc}
\usepackage{cite}
%\usepackage{listings}
\usepackage[usenames]{xcolor}
\usepackage{graphicx}
\usepackage{hyperref}
\usepackage{xr} % cross-reference across files
\usepackage{float} %for H positioning
\usepackage{listings}
\usepackage{makecell}
\usepackage{gitinfo2}

\externaldocument[TUT-]{tutorial}

% start a new section on a new page
\let\stdsection\section
\renewcommand\section{\newpage\stdsection}

%\lstset{ %
%  basicstyle=\footnotesize,        % the size of the fonts that are used for the code
%  breakatwhitespace=true,          % sets if automatic breaks should only happen at whitespace
%  breaklines=true,                 % sets automatic line breaking
%  commentstyle=\color{gray},    % comment style
%  keepspaces=true,                 % keeps spaces in text, useful for keeping indentation of code (possibly needs columns=flexible)
%  language=bash                 % the language of the code
%}

%opening
\title{CTA user manual}
%\author{Roger Biffiger, PSI}
\date{}

\begin{document}

\maketitle

\version{This user manual has been generated from the CTA src tree version \gitFirstTagDescribe}

%\begin{abstract}
%\end{abstract}

\tableofcontents

%\newpage

\section{Introduction}\label{sec:Introduction}
This is the user manual for the CTA (Complex Triggering Application).
The CTA allows a user to define a sequence of events and to distribute
these events to the event receivers in a subtree of a MRF (Micro
Research Finland) based timing system. In SwissFEL CTA is used in the
experimental stations ESA and ESB.

\section{Overview}\label{sec:Overview}
A timing system build from MRF hardware components has a tree topology.
There is one master at the root of the tree which generates the timing
stream signal. There are fanouts at each
crotch to multiply the number of ports providing a timing stream signal.
Finally there are event receivers at the leafs of the tree which receive
the timing stream, decode it and produce pulses on physical outputs
to trigger devices.
A simplified timing system tree can be seen in Figure \ref{fig:tim_sys_tree}.

\begin{figure}[H]
	\centering
	\includegraphics[width=\columnwidth]{./img/overview_1}
	\caption{Simplified timing system tree}
	\label{fig:tim_sys_tree}
\end{figure}

The decision which event code is distributed with what delay in each machine
pulse is done by the two EPICS software applications TMA (Timing Master
Application) and CTA (Complex Triggering Application).

The TMA runs on the master and has the following responsibilities:
\begin{itemize}
\item
Provide a set of PVs to the operator which allows the operator to control
the timing of the accelerator.
\item
Based on the operator PVs, the TMA must compute the sequence of events which
is to be send out in the next machine pulse.
\item
The TMA must write the event sequence to the hardware which writes it
to the timing stream.
\end{itemize}

The CTA is an EPICS software application which runs on a fanout. It has the
following responsibilities:
\begin{itemize}
\item
Provide a user interace which allows the user to define a sequence of event codes.
\item
The user must be able to download the event code sequence to the IOC
running the CTA.
\item
The CTA must program the hardware such that the sequence of event codes
is send in the timing stream.
\end{itemize}

TMA event codes are seen by all event receivers in the tree. CTA event
codes are seen by event receivers in the CTA subtree only

\section{CTA overview}\label{sec:CTA}

\subsection{Concept and terminology}

\begin{figure}[H]
	\centering
	\includegraphics[width=\columnwidth]{./img/cta_overview_1}
	\caption{CTA overview}
	\label{fig:cta_overview}
\end{figure}

Figure \ref{fig:cta_overview} shows the context in which the CTA is running.
As shown in the figure, there are three different user interfaces to control the
CTA. These user interfaces are described in \ref{sec:CtaUserInterfaces}.

As mentioned above the user can define a sequence of events.
A sequence consists of one or more series and a series defines, if a specific event
shall be sent in a machine pulse or not.

\begin{table}[!hbt]
\caption{CTA sequece}
\label{tab:cta_sequence}
\centering
  \begin{tabular}{|c|c|c|}
    \hline \textbf{machine pulse} & \textbf{event code A} & \textbf{event code B}\\ \hline
    \hline X     & 1 & 0\\
    \hline X + 1 & 0 & 1\\
    \hline X + 2 & 1 & 0\\
    \hline
  \end{tabular}
\end{table}

Table \ref{tab:cta_sequence} shows a CTA sequence. In this sequence event code A is
send in machine pulses X and X + 2. Event code B is send in machine pulse X + 1.

Using the user interface the user can do the following actions:
\begin{itemize}
\item
The sequence can be uploaded (user interface to CTA).
\item
The sequence can be downloaded (CTA to user interface).
\item
The number of times the sequence is repeated by the CTA can be set (integer or
forever).
\item
The start condition can be set. When the condition is met, the CTA starts to
send out the sequence of event codes in the timing stream. Currently the
options to start the sequence immediately or based on the pulse id are
supported.
\item
The CTA can be instructed to start inserting the sequence in the timing stream.
\item
The CTA can be instructed to stop inserting the sequence in the timing stream.
\item
The status of the CTA can be retrieved. It shows whether the CTA is running
or stopped and it shows the pulse id at which the last sequence was started.
\end{itemize}

Refer to section \ref{sec:CtaExampes} for examples how an event code sequence
is injected in a timing stream signal.

\subsection{Event and delay space}
To avoid interference of TMA and CTA event codes, TMA and CTA event codes
are separated in event code space and time.

\begin{itemize}
\item
Event codes from 200-219 are reserved for CTA events. The TMA will not
use these event codes.
\item
The delay of CTA event codes is between 100-233 us. The TMA will not send event
codes with this delay.
\end{itemize}
ESA CTA event codes can not interfere with ESB CTA events codes because the
subtrees are on separate branches of the timing tree.

CTA event codes have a fixed delay. The delay value for each CTA event code
can be seen in table \ref{tab:cta_event_delays}

\begin{table}[!hbt]
\caption{Delay of CTA event codes}
\label{tab:cta_event_delays}
\centering
  \begin{tabular}{|c|c|}
    \hline \textbf{event code} & \textbf{delay [us]} \\ \hline
    \hline 200 & 100\\
    \hline 201 & 107\\
    \hline 202 & 114\\
    \hline 203 & 121\\
    \hline 204 & 128\\
    \hline 205 & 135\\
    \hline 206 & 142\\
    \hline 207 & 149\\
    \hline 208 & 156\\
    \hline 209 & 163\\
    \hline 210 & 170\\
    \hline 211 & 177\\
    \hline 212 & 184\\
    \hline 213 & 191\\
    \hline 214 & 198\\
    \hline 215 & 205\\
    \hline 216 & 212\\
    \hline 217 & 219\\
    \hline 218 & 226\\
    \hline 219 & 233\\
    \hline
  \end{tabular}
\end{table}

\section{CTA user interfaces}\label{sec:CtaUserInterfaces}
\subsection{GUI}
\subsubsection{Starting the GUI}
The GUI to control the CTA can be started on a console with the following
command:
\begin{itemize}
\item
ESA: start\_cta\_gui.sh ESA SAR-CCTA-ESA
\item
ESB: start\_cta\_gui.sh ESB SAR-CCTA-ESB
\end{itemize}

\subsubsection{Using the GUI}

The CTA GUI is shown in figure \ref{fig:cta_gui}.

\begin{figure}[H]
	\centering
	\includegraphics[]{./img/gui_1}
	\caption{CTA GUI}
	\label{fig:cta_gui}
\end{figure}

\paragraph{sequence upload/download:}
The widgets in the section \emph{sequence upload/download} can be used
to upload and download the sequence.
A widget indicates whether the sequence in the sequence definition table
and the sequence on the IOC are equal or not.

\paragraph{sequence definition:}
The section \emph{sequence definition} contains the sequence definition
table. Each line in the table defines an event code which is to be distributed
in a certain machine pulse.
The \emph{Step Offset} and \emph{Start Offset} columns define
the machine pulse. The column \emph{Event Code} defines the event code.
The \emph{Step Offset} can be used to define the machine pulse for sending
relative to the last line and the \emph{Start Offset} can be used to define
the machine pulse for sending relative to the first line.
If the user enters a \emph{Step Offset} the \emph{Start Offset} is
set automatically and vice versa.
There are buttons to insert and remove new lines. This can also be
done with the context menu.

\paragraph{start configuration:}
The radio buttons in the \emph{start configuration} section can be used to
define the mode, which is used by the CTA, to starts sending out the sequence on
the timing stream.
Currently the modes \emph{start immediately} and
\emph{start if ((pulseId \% divisor) - offset == 0)} are supported.
In the mode \emph{start immediately} the CTA starts sending the event sequence
immediately. The time between pressing the start button and the time the first
event code is send is short but undefined.
In the mode \emph{start if ((pulseId \% divisor) - offset == 0)} the first
event code of the sequence is sent out when the if condition is met for the
first time.

\paragraph{repetition configuration:}
The radio buttons in the \emph{repetition configuration} section can be used
to define how many times the sequence is repeated.

\paragraph{run status:}
The widgets in the \emph{run status} section can be used to start or stop
event code sending of the CTA. Additionally the \emph{status} widget shows
the current status of the CTA (running or stopped) and the \emph{started at}
widget shows the pulse id at which the last sequence has been started.

\subsection{Channel access}

Table \ref{tab:cta_public_pvs} lists the public PVs of CTA. Note, that currently
only one sequence is supported. \$(P) is the device name of CTA (e.g. SAR-CCTA-ESA
for the CTA in ESA).

\begin{table}[!hbt]
\caption{CTA public PVs}
\label{tab:cta_public_pvs}
\centering
  \begin{tabular}{|c|c|}
    \hline \textbf{PV name} & \textbf{comment} \\ \hline
    \hline \$(P):seq0Ser0-Data-I & sequence 0 series 0 (event code 200)\\
    \hline \$(P):seq0Ser1-Data-I & sequence 0 series 1 (event code 201)\\
    \hline : & :\\
    \hline \$(P):seq0Ser19-Data-I & sequence 0 series 19 (event code 219)\\
    \hline \$(P):seq0Ctrl-Length-I & length of sequence\\
    \hline \$(P):seq0Ctrl-Cycles-I & number of repetitions (0 = forever, n = n times)\\
    \hline \$(P):seq0Ctrl-Start-I & start sequece (caput pv 1)\\
    \hline \$(P):seq0Ctrl-Stop-I& stop sequence (caput pv 1)\\
    \hline \$(P):seq0Ctrl-SCfgMode-I &
           \makecell {start config mode\\
                      (0 = immediately,\\
                      1 = if ((pulse id \% divisor) - offset == 0))}\\
    \hline \$(P):seq0Ctrl-SCfgModDivisor-I & start config modulo divisor\\
    \hline \$(P):seq0Ctrl-SCfgModOffset-I & start config modulo offset\\
    \hline \$(P):seq0Ctrl-IsRunning-O & 0 = stopped, 1 = running\\
    \hline \$(P):seq0Ctrl-StartedAt-O & pulse id when sequence was started\\
    \hline \$(P):SerMaxLen-O & max length of sequence\\
    \hline
  \end{tabular}
\end{table}

The length of a sequence is the number of machine pulses for which it defines the
CTA event codes.

Listing \ref{lst:caput} show an example of how to upload and start a simple sequence
on CTA in ESA.  

\begin{minipage}{\linewidth}
\begin{lstlisting}[caption={caput listing}, label={lst:caput}, linewidth=\columnwidth]
caput SAR-CCTA-ESA:seq0Ser0-Data-I "1 0 1"
caput SAR-CCTA-ESA:seq0Ser1-Data-I "0 1 0"
caput SAR-CCTA-ESA:seq0Ser2-Data-I "0 0 0"
caput SAR-CCTA-ESA:seq0Ser3-Data-I "0 0 0"
caput SAR-CCTA-ESA:seq0Ser4-Data-I "0 0 0"
caput SAR-CCTA-ESA:seq0Ser5-Data-I "0 0 0"
caput SAR-CCTA-ESA:seq0Ser6-Data-I "0 0 0"
caput SAR-CCTA-ESA:seq0Ser7-Data-I "0 0 0"
caput SAR-CCTA-ESA:seq0Ser8-Data-I "0 0 0"
caput SAR-CCTA-ESA:seq0Ser9-Data-I "0 0 0"
caput SAR-CCTA-ESA:seq0Ser10-Data-I "0 0 0"
caput SAR-CCTA-ESA:seq0Ser11-Data-I "0 0 0"
caput SAR-CCTA-ESA:seq0Ser12-Data-I "0 0 0"
caput SAR-CCTA-ESA:seq0Ser13-Data-I "0 0 0"
caput SAR-CCTA-ESA:seq0Ser14-Data-I "0 0 0"
caput SAR-CCTA-ESA:seq0Ser15-Data-I "0 0 0"
caput SAR-CCTA-ESA:seq0Ser16-Data-I "0 0 0"
caput SAR-CCTA-ESA:seq0Ser17-Data-I "0 0 0"
caput SAR-CCTA-ESA:seq0Ser18-Data-I "0 0 0"
caput SAR-CCTA-ESA:seq0Ser19-Data-I "0 0 0"
caput SAR-CCTA-ESA:seq0Ctrl-Length-I 3
caput SAR-CCTA-ESA:seq0Ctrl-Cycles-I 1
caput SAR-CCTA-ESA:seq0Ctrl-Start-I 1
\end{lstlisting}
\end{minipage}

\subsection{Python library}
The python library is available in the PSI python 3.6 environment. You can get
this environment with the following command: source /opt/gfa/python 3.6
Use the following commands to create a CTA library object, upload a sequence
to the ESA CTA and start it.

\begin{lstlisting}[caption={python CTA lib}, label={lst:python_cta_lib}, language=Python]
import cta_lib

# create cta lib object
lib = cta_lib.CtaLib("SAR-CCTA-ESA")

# create sequence
sequence = {}
sequence[200] = [1, 0, 1]
sequence[201] = [0, 1, 0]

# upload
lib.upload(sequence)

# start
lib.start()
\end{lstlisting}

Refer to the python docstrings for more information about the available
operations.
Furthermore a few example scripts which demonstrate the usage of the python
library are available in the PSI python environment.
\begin{itemize}
\item
cta\_lib\_ex01\_upload\_download\_print
\item
cta\_lib\_ex02\_upload\_repeat\_n\_wait
\item
cta\_lib\_ex03\_upload\_run\_forever\_stop
\item
cta\_lib\_ex04\_callbacks
\item
cta\_lib\_ex05\_start\_modulo.py
\end{itemize}
You can call these scripts from the terminal command line. Use the -h option to get help.

\section{CTA usage examples}\label{sec:CtaExampes}
This section provides some examples of CTA GUI configurations together
with the resulting sequence of event codes.

\subsection{One local sequence started immediately}\label{subsec:one_seq_imm}
\begin{figure}[H]
	\centering
	\includegraphics[width=\columnwidth]{./img/ex_single_seq_1}
	\caption{GUI and resulting event code sequence}
	\label{fig:pulser_generic}
\end{figure}

\subsection{One local sequence started if modulo expression is true}\label{subsec:one_seq_modulo}
\begin{figure}[H]
	\centering
	\includegraphics[width=\columnwidth]{./img/ex_single_seq_2}
	\caption{GUI and resulting event code sequence}
	\label{fig:pulser_generic}
\end{figure}

\subsection{Generating 50 Hz event with CTA}\label{subsec:one_seq_50Hz}
\begin{figure}[H]
	\centering
	\includegraphics[width=\columnwidth]{./img/50_hz}
	\caption{GUI and resulting event code sequence}
\end{figure}

\end{document}
